%Перечень условных обозначений (при необходимости) оформляется в алфавитном порядке
\sectioncentered*{Перечень условных обозначений}
\addcontentsline{toc}{section}{Перечень условных обозначений}
\label{sec:reduction}
\begin{enumerate}
    \item API -- Application Programming Interface (интерфейс программирования приложений);
    \item ASGI -- Asynchronous Server Gateway Interface (асинхронный шлюзовой интерфейс сервера);
    \item BFS/DFS -- Breadth-First/Depth-First Search (алгоритмы обхода графа в ширину/в глубину);
    \item БЗ -- База знаний;
    \item CI/CD -- Continuous Integration/Continuous Delivery (непрерывная интеграция и доставка);
    \item CLI -- Command Line Interface (интерфейс командной строки);
    \item CS -- Computer Science (фундаментальные основы информатики);
    \item EdTech -- Educational Technology (образовательные технологии);
    \item FAANG -- Facebook, Amazon, Apple, Netflix, Google (ведущие технологические компании);
    \item Git -- распределённая система контроля версий;
    \item HR-tech -- Human Resources Technology (технологии управления персоналом);
    \item JSON -- JavaScript Object Notation (формат обмена данными);
    \item LRU -- Least Recently Used (алгоритм вытеснения наименее недавно использованных данных);
    \item MVP -- Minimum Viable Product (минимально жизнеспособный продукт);
    \item OOP -- Object-Oriented Programming (объектно-ориентированное программирование);
    \item PR -- Pull Request (запрос на слияние изменений в репозиторий);
    \item RAG -- Retrieval-Augmented Generation (генерация с дополненным поиском);
    \item REST -- Representational State Transfer (архитектурный стиль проектирования API);
    \item SLA/SLO/SLI -- Service Level Agreement/Objective/Indicator (соглашение/цель/индикатор уровня обслуживания);
    \item SPA -- Single Page Application (одностраничное приложение);
    \item STAR -- Situation, Task, Action, Result (методика структурирования поведенческих ответов);
    \item TechInterview -- интеллектуальный помощник по подготовке к техническим собеседованиям;
    \item TTL -- Time To Live (время жизни записи или сессии);
    \item UI/UX -- User Interface/User Experience (пользовательский интерфейс/пользовательский опыт);
\end{enumerate}


